\documentclass[11pt]{article}
\usepackage[a4paper,margin=28mm]{geometry}
\usepackage{amsmath,amssymb,amsthm}
\usepackage[hidelinks]{hyperref}
\newtheorem{theorem}{Theorem}
\newtheorem{lemma}[theorem]{Lemma}
\newtheorem{corollary}[theorem]{Corollary}
\title{Polynomial thresholds for graceful leaf extensions of arbitrary trees}
\author{Aky-lab}\date{Working preprint, 25 November 2026}
\newcommand{\proofmodule}[2]{\href{https://github.com/Aky-lab/graceful-leaf-threshold-preprint/blob/v1.1/lean/Next/#1.lean}{#2}}
\begin{document}
\maketitle
\begin{abstract}
Let $H$ be any tree on $n$ vertices, and attach new leaves at its vertices
in any prescribed numbers. The resulting tree is graceful whenever its
final edge count is at least $3{,}955{,}201 n^7$. No graceful labeling of
$H$ is assumed, and no positive lower bound is required at an individual
attachment site. The proof uses the parities of the same prescribed leaf counts to
control two competing constructions. For a tree with $h$ internal
vertices the resulting sufficient edge bound is $3{,}955{,}201h^7$.
The construction lemmas and the uniform estimate are formalized in Lean.
\end{abstract}
\section{The threshold theorem}
A tree with $N$ edges is \emph{graceful} if it admits a bijection from
its vertices to $\{0,\ldots,N\}$ whose edge differences are exactly
$1,\ldots,N$. For a finite tree $H$ and a vector
$\mathbf a\in\mathbb Z_{\ge0}^{V(H)}$, let $H(\mathbf a)$ be obtained
by adding $a_v$ new leaves at each $v\in V(H)$. These are the only new
vertices; the construction must realize the entire prescribed target.
\begin{theorem}\label{poly}
Let $H$ be a tree on $n\ge1$ vertices. Every prescribed leaf extension
$H(\mathbf a)$ with final edge count
\[
 N=n-1+\sum_{v\in V(H)}a_v\ \ge\ 3{,}955{,}201n^7
\]
is graceful.
\end{theorem}
The hypothesis uses the number of vertices before attachment and the number
of edges after attachment. It allows all new leaves at a single vertex.
The general Graceful Tree Conjecture remains open.

Superdock proved eventual gracefulness with all new leaves at a single
prescribed vertex \cite[Corollary 5.10]{S}. Iterating such a result lets
later thresholds depend on earlier leaf counts; it does not give the
uniform bound above. Lapierre and McGuinness \cite[Theorem 1.2 and
Corollary 1.3]{LM} give constructions from bipartite gracefully labeled
cores with conditions on the attachment demands. Here no graceful
source labeling or bipartition-demand hypothesis is imposed.
We use the classical pair-transfer identity, together with a modular
construction controlling several attachment ranks and a two-interval
allocation for all secondary sites.

\subsection{Proof strategy and notation}
Call an injective nonnegative rooted labeling \emph{consistent} if its
root is labeled zero and its positive vertex labels, as a multiset,
are exactly its edge differences. We call labels \emph{ranks}. Scaling
preserves consistency. If a consistent tree has ranks in $[0,N]$, an
unused positive rank may be used for a new root leaf: its label and its
edge difference coincide. Thus the task is to meet all nonroot demands
before filling the complement at zero.

Two elementary operations preserve this equality of multisets. At a
parent of rank $p$, new children of ranks $x,p-x$, with $0<x<p$,
supply differences $p-x,x$. Also leaves of ranks $x,y$ can be moved
between parents of ranks $p,q$ when $x+y=p+q$ and the corresponding
absolute differences are exchanged. The two constructions below use only vertices of the prescribed extension.
We state their quantitative guarantees explicitly. The linked Lean modules
give the detailed labeling constructions.

\section{Completion from one or two large supplies}
Choose a vertex $u$ of maximum demand, let $q$ count vertices of positive
demand, and let $o$ count vertices other than $u$ whose demand is odd.
Under the edge hypothesis of Theorem~\ref{poly}, some demand is positive,
so $0\le o<q\le n$. Write $S=n+o$. For $q\ge2$ put
\[
 T=(q-1)^2\bigl(1236n+1200q-3764-1200\max(0,o-1)\bigr),
 \qquad L=(q-1)(T-2).
\]
For $q=1$ set $L=0$; the two-site case is then absent.
Define the seed and packing allowances by
\[
 A(S)=4096S^6(9+3\lceil\log_2 S\rceil),\qquad
 P(S,L)=200S^3\bigl(2(S-1)+2L+1\bigr).
\]

\begin{lemma}[Completion alternatives]\label{routes}
If $q\ge2$ and some $v\ne u$ has demand at least $T$, the extension is graceful.
In the absence of such a site, the extension is graceful whenever
$N\ge A(S)+P(S,L)$.
\end{lemma}
\begin{proof}
In the first case both $u$ and $v$ meet the threshold, by maximality of $u$.
For each of the other positive sites reserve one requested leaf if its
demand is odd and two if it is even. If $r$ is the number reserved, then
\[
 r+o=2(q-2)+(a_v\bmod 2).
\]
Thus the enlarged core has at most $j=n+2q-5-\max(0,o-1)$ edges.
The \proofmodule{ParityPacketTwoBank}{two-site construction} uses a
modular labeling and transfers leaves between the two chosen sites.
For $m=q-1$ its sufficient demand at each is
$40m^2(30j+25-30m)+36m^2(n+1)$, which expands to $T$.
The \proofmodule{ParityPacketPalette}{palette estimates} give enough
available ranks in the two intervals, and the construction also handles
its exceptional double-broom case.

In the other case include one requested leaf at each odd site off $u$.
The prepared core has order $S$ and every remaining demand off $u$ is
even. When $q\ge2$, $T$ is even, so each such residual demand is at most
$T-2$. Their total is at most $L$ (zero also when $q=1$).
The \proofmodule{SparsePackingAssembly}{one-site packing bound}, applied in
\proofmodule{ParityPacketAssembly}{the completion lemma}, uses $A(S)$ for
the prepared rooted labeling and $P(S,L)$ to make room for the residual
pairs. At rank $p$ these pairs have the form $\{x,p-x\}$; unused ranks
then complete the demand at the root $u$. This gives a labeling of the
entire prescribed extension.
\end{proof}

For example, consider the six-vertex path $0,1,2,3,4$ with an extra
leaf $5$ at vertex $1$. Prescribe respectively $24{,}209$, $24{,}209$,
and one new leaf at $0,4,5$. With $u=0$ we have $q=3$ and $o=2$.
The threshold at each of $0$ and $4$ is $24{,}208$, so the first case
already constructs the target. Ignoring the two off-root parity bits
would give a threshold of $29{,}008$. This example is included in the
\proofmodule{ParityPacketRegression}{checked regression examples}.

\section{The uniform estimate}
\begin{proof}[Proof of Theorem~\ref{poly}]
By Lemma~\ref{routes} it suffices to show
$A(S)+P(S,L)\le3{,}955{,}201n^7$. The correlation between $S$ and
$T$ matters: an additional odd site enlarges the prepared core but also
reduces the two-site reserve. For $n\ge6$ write $B=2n-1$ and
$C=3636n-2564$. Then $S\le B$. Put
\[
 U=\max(1,q-1)^2(1236n+1200q-1200o-2564).
\]
We have $L\le(q-1)U$ by the formula for $T$ when $q\ge2$, and
$2(S-1)+1\le U$ since $o<q$ and $S\le2n-1$.
Consequently
\[
 2(S-1)+2L+1\le B U,\qquad U\le(n-1)^2(C-1200S).
\]
The function $s^3(C-1200s)$ is increasing on $[0,B]$, since
its derivative is $3s^2(C-1600s)$ and $1600B\le C$.
Hence the packing cost is at most
\[
 E(n)=200(2n-1)^4(n-1)^2(1236n-1364).
\]
The leading coefficient here is $200\cdot 2^4\cdot1236=3{,}955{,}200$.

For the seed, set $t=\lfloor S/2^{21}\rfloor$. Since
$S\le2^{21}(t+1)\le2^{21+t}$ and $S\le2n$, we get
\[
 A(S)\le4096(2n)^6(72+3t)\le18{,}874{,}368n^6+n^7.
\]
For $n\ge32$, substituting $n=x+32$ into
$3{,}955{,}200n^7-E(n)-18{,}874{,}368n^6$ gives a polynomial with
nonnegative coefficients. In ascending order of degree these are
\[
\begin{gathered}
10{,}901{,}718{,}756{,}768,\quad 44{,}127{,}171{,}891{,}744,\quad
6{,}783{,}308{,}678{,}880,\\
423{,}670{,}516{,}320,\quad 13{,}283{,}768{,}320,\quad
208{,}588{,}544,\quad 1{,}311{,}232 .
\end{gathered}
\]
This absorbs the degree-six seed term and
leaves one further copy of $n^7$ for the seed. For $6\le n<32$ direct
integer substitution gives $A(2n-1)+E(n)\le3{,}955{,}201n^7$.
For $1\le n\le5$ substitute the finitely many $1\le q\le n$ and
$0\le o<q$ directly into $A(S)+P(S,L)$. The
\proofmodule{ParityPacketUniform}{formal uniform estimate} includes
the polynomial identity and the exact small-case checks. This proves the
bound in both completion cases.
\end{proof}

\section{Effective consequences}
\begin{corollary}\label{effective}
For each finite tree $H$ on $n$ vertices, the least nonnegative integer
$B(H)$ such that every extension with at least $B(H)$ new leaves is
graceful is computable, and
\[
 B(H)\le 3{,}955{,}201n^7-(n-1).
\]
The complete finite list of nongraceful attachment vectors is computable.
\end{corollary}
\begin{proof}
Set $U=3{,}955{,}201n^7-(n-1)$. Enumerate all nonnegative integer
vectors on $V(H)$ of total less than $U$. Each resulting finite tree
can be tested for gracefulness by enumerating its vertex bijections.
Theorem~\ref{poly} certifies every omitted vector. If no vector fails,
return $B(H)=0$; otherwise return one plus the largest failing total.
This proves computability, not practical efficiency of this enumeration.
\end{proof}
\begin{samepage}
\begin{corollary}\label{internalorder}
A nongraceful tree with $N$ edges and $h$ internal vertices must satisfy
\[
 N<3{,}955{,}201h^7,\qquad
 h>(N/3{,}955{,}201)^{1/7}.
\]
For each fixed $h$, all nongraceful trees with at most $h$ internal
vertices can therefore be listed by a finite computation.
\end{corollary}
\begin{proof}
Trees of order at most two are graceful. In any larger tree the internal
vertices induce a nonempty subtree, and all other vertices are leaves
adjacent to it. Apply Theorem~\ref{poly} to that core. For bounded $h$
only finitely many trees remain below the displayed edge bound.
\end{proof}
\end{samepage}

\section*{Checking the formal proof}
Clone the \href{https://github.com/Aky-lab/graceful-leaf-threshold-preprint/tree/v1.1}{tagged public source}
and enter its \texttt{lean/} directory. With Elan installed, run
\texttt{lake exe cache get} followed by \texttt{lake build}.
The manifest pins Mathlib and the toolchain file selects Lean 4.30.0.
The default target compiles the \proofmodule{TreeEncodingPolynomialParityPacket}{graph-level theorems}
for the prescribed extension and for the internal-vertex bound.
It checks their axiom dependencies against propositional extensionality,
classical choice, and quotient soundness.
\small\begin{thebibliography}{9}
\bibitem{LM} L. Lapierre and S. McGuinness, Constructing trees with graceful labelings using caterpillars, \emph{J. Combin. Math. Combin. Comput.} \textbf{94} (2015), 313--321.
\bibitem{S} M. Superdock, The Graceful Tree Conjecture: a class of graceful diameter-6 trees, \href{https://arxiv.org/abs/1403.1564}{arXiv:1403.1564} (2014).
\end{thebibliography}
\end{document}
